% ============================================
% 2. 问题分析
% 小标题格式：2.N 问题N——基于XX的XX分析
% ============================================

\section{问题分析}

\subsection{问题一——基于【主要方法】的【问题类型】分析}

针对问题一，该问题的本质是一个\textbf{【问题类型：预测/分类/评价/优化】}问题。考虑到附件数据具有\textbf{【数据特征：高维度/时序性/非线性/多变量】}的特点，拟采用\textbf{【主要模型/方法】}进行求解。

分析流程如下：首先，对附件数据进行\textbf{【预处理方式】}，提取\textbf{【关键特征】}；接着，建立\textbf{【模型名称】}模型，通过\textbf{【参数确定方法】}确定模型参数；然后，利用\textbf{【求解算法】}进行求解；最后，对模型结果进行\textbf{【验证方式】}和\textbf{【灵敏度分析】}。

具体分析流程如图\ref{fig:analysis1_flow}所示。

\begin{figure}[ht]
  \begin{center}
    \begin{tikzpicture}[x=1cm, y=0.7cm]
      \node[box] (n11) at (0,0) {数据预处理};
      \node[box] (n12) at (5,0) {特征提取};
      \node[box] (n13) at (10,0) {模型构建};
      \node[box] (n22) at (5,-3) {模型求解};
      \node[box] (n21) at (0,-3) {结果验证};
      \node[box] (n23) at (10,-3) {参数优化};
      \node[box] (n31) at (0,-6) {结果输出};
      \node[box] (n32) at (5,-6) {灵敏度分析};
      \node[box] (n33) at (10,-6) {结论总结};
      \draw[arrow] (n11)--(n12); \draw[arrow] (n12)--(n13);
      \draw[arrow] (n23)--(n22); \draw[arrow] (n22)--(n21);
      \draw[arrow] (n13) -- ++(0,-1.0) -| (n23);
      \draw[arrow] (n21)--(n31); \draw[arrow] (n31)--(n32); \draw[arrow] (n32)--(n33);
    \end{tikzpicture}
    \caption{问题一分析流程图}
    \label{fig:analysis1_flow}
  \end{center}
\end{figure}

\subsection{问题二——基于【主要方法】的【问题类型】分析}

针对问题二，该问题的本质是一个\textbf{【问题类型】}问题。考虑到\textbf{【问题特点】}，拟采用\textbf{【主要模型/方法】}进行求解。

分析流程如下：首先，\textbf{【步骤1】}；接着，\textbf{【步骤2】}；然后，\textbf{【步骤3】}；最后，\textbf{【步骤4】}。

具体分析流程如图\ref{fig:analysis2_flow}所示。

\begin{figure}[ht]
  \begin{center}
    \begin{tikzpicture}[x=1cm, y=0.7cm]
      \node[box] (n11) at (0,0) {...};
      \node[box] (n12) at (5,0) {...};
      \node[box] (n13) at (10,0) {...};
      \node[box] (n22) at (5,-3) {...};
      \node[box] (n21) at (0,-3) {...};
      \node[box] (n23) at (10,-3) {...};
      \node[box] (n31) at (0,-6) {...};
      \node[box] (n32) at (5,-6) {...};
      \node[box] (n33) at (10,-6) {...};
      \draw[arrow] (n11)--(n12); \draw[arrow] (n12)--(n13);
      \draw[arrow] (n23)--(n22); \draw[arrow] (n22)--(n21);
      \draw[arrow] (n13) -- ++(0,-1.0) -| (n23);
      \draw[arrow] (n21)--(n31); \draw[arrow] (n31)--(n32); \draw[arrow] (n32)--(n33);
    \end{tikzpicture}
    \caption{问题二分析流程图}
    \label{fig:analysis2_flow}
  \end{center}
\end{figure}

\subsection{问题三——基于【主要方法】的【问题类型】分析}

针对问题三，该问题的本质是一个\textbf{【问题类型】}问题。考虑到\textbf{【与前序问题的递进关系】}，拟采用\textbf{【主要模型/方法】}进行求解。

分析流程如下：首先，\textbf{【步骤1】}；接着，\textbf{【步骤2】}；然后，\textbf{【步骤3】}；最后，\textbf{【步骤4】}。

\subsection{问题四——基于【主要方法】的【问题类型】分析}

针对问题四，该问题的本质是一个\textbf{【问题类型】}问题。基于\textbf{【前三个问题的分析结果】}，拟从\textbf{【角度1】}和\textbf{【角度2】}出发，给出\textbf{【建议/改进方案/推广方向】}。
